\section{Discussion and Open Questions}
\label{sec:discussion}

\subsection{What this paper establishes}

\paragraph{Need-sufficiency resolves the polarity problem.}
The sufficiency threshold identifies the structural boundary
where sacrifice polarity flips.  Below-sufficiency sacrifices
are involuntary and correctly filtered by S1
(Proposition~\ref{prop:s1_sufficiency}).  The need-access cost
diagnostic (Definitions~\ref{def:nac}--\ref{def:nse}) tracks
need-satisfaction efficiency as a companion metric.

\paragraph{The B-to-C gap decomposes into disjoint
diagnostic cells.}
Under poset-boundary closure, $1 - \betaL =
\delta_{\mathrm{restricted}} + \delta_{\mathrm{covered}} +
\delta_{\mathrm{dormant}} + \delta_{\mathrm{residual}}$
(Definition~\ref{def:gap_decomp}).  The partition is
computable in $O(|P| + E + M)$
(Proposition~\ref{prop:gap_computable}) given the capability
census, poset structure, restriction and residual sets, event
admissibility labelling, and a closure policy.  Alarm firings
can be interpreted diagnostically by which $\delta$ dominates.

\paragraph{A wireheading-consistent signal is computable from
trade data.}
Trade-flow concentration
(Proposition~\ref{prop:wireheading_hhi}) provides a tractable
concentration signal---Schur-convex in the majorization order
and monotone in the concentration of observed trade flow---that
is \emph{consistent with} wireheading but does not constitute a
wireheading determination on its own.  The Part-A category-flow
variant is third-party observable from public committed
sacrifice events
(Proposition~\ref{prop:third_party_wireheading}).  Interpreting
an $\HHI_{\mathrm{alarm}}$ crossing as wireheading requires
additional context (governance-set threshold, domain baseline
concentration, and cross-diagnostic confirmation via the gap
decomposition and need-sufficiency metrics).

\paragraph{The residual class is structurally distinct.}
Capabilities in $\ResS$
(Definition~\ref{def:residual_class}) have zero sacrifice
evidence by construction and cannot gain it regardless of
observation time.  When $\ResS$ is poset-isolated from the
non-residual cells (the residual-cell irreducibility hygiene
check of Remark~\ref{rem:residual_cell_irreducibility}), the
residual class is the irreducible component of the gap
decomposition under priority-extension closure; without
isolation, residual mass can migrate to neighbouring cells
under priority-extension and the boundary-residual or
sub-partition variants of
Remark~\ref{rem:cell_cooperative_closure} are needed to
preserve attribution.  The revealed-observation carve of
$\ResS$ (as the set of capabilities \emph{structurally outside
the channel's reach}) differs from and is operationally
cleaner than the classical benchmarkability-based carve.

\subsection{Limitations}

\paragraph{Threshold governance is not framework-derived.}
The sufficiency thresholds $s_k$ and the admissible
reference set $\Pi_N^{\mathrm{ref}}$
(Definition~\ref{def:downstream_safe}) are governance inputs,
not framework-computed.  The framework reads them at evaluation
time and does not attempt to derive them from first
principles.

\paragraph{The filter's interaction with capability-stuffing.}
Proposition~\ref{prop:below_suff_saturation} establishes a
pointwise below-sufficiency cap.  This is \emph{not} an
aggregate dimension-level cap: distinct capabilities each
satisfying sufficiency on the same dimension each score $1$
and accumulate under M5.  Full structural resolution of
capability-stuffing requires machinery outside this paper's
scope.

\paragraph{The gap decomposition requires poset-boundary
closure.}
The exact identity $1 - \betaL = \sum \delta$ holds only when
the cells $\{R, C, D, S\}$ are poset-boundary closed (CC1 +
CC2 in Definition~\ref{def:gap_decomp}).
Remark~\ref{rem:cell_cooperative_closure} describes three
strategies to achieve closure in practice.

\paragraph{The residual class is characterized but not valued.}
The residual class (Section~\ref{sec:residual_class})
characterizes \emph{which} capabilities the framework cannot
reach, not \emph{how} to value them.

\subsection{Open questions}

\begin{enumerate}

\item \label{oq:threshold_governance}
\textbf{Threshold governance.}
Several interrelated questions surround the sufficiency
thresholds.  (a)~\emph{Sensitivity}: if the threshold is set
too low, need-costs between the threshold and true sufficiency
are misclassified as want-pursuit events; if too high, the
framework over-constrains the need.
(b)~\emph{Dynamic update}: thresholds may shift as technology
and social context evolve; a mechanism for threshold revision
(analogous to controlled relaxation applied to sufficiency
standards) is needed.
(c)~\emph{Cross-cultural variation}: different populations may
set thresholds differently; can the framework maintain
culturally-indexed thresholds while preserving a coherent
aggregate diagnostic, or does variation fragment it?

\item \label{oq:need_identification}
\textbf{Need identification and M5 interaction.}
How does the framework identify which capabilities are needs
versus wants?  The poset-structural approach (high fan-out $=$
need) is a heuristic, not infallible; a hybrid of structural
position and revealed-sacrifice polarity transition is likely
required.  Relatedly, under the capped-benchmark discipline
(Definition~\ref{def:polarity_benchmark}), sufficiency-met
needs contribute a fixed cap whose weight is a design
parameter affecting the relative importance of
need-satisfaction versus want-pursuit in $\volL$.

\item \label{oq:capability_stuffing}
\textbf{Capability-stuffing as structural avoidance.}
Capability-stuffing (Section~\ref{sec:capability_stuffing})
remains an open M5 pathology under this paper's machinery.  No
signature over post-declaration poset structure or exercise
pattern separates stuffed from legitimate capabilities in a way
the agent cannot adversarially manipulate.  The canonical
example reveals itself as an instance of
\emph{structural avoidance} flagged as an open gap
in~\citep[Discussion]{lasser2026poset}: the agent declines to
discover that ``travel from home to gas station'' and ``travel
from work to gas station'' admit merger into a canonical
multi-origin fuel-access capability.  Related avoidance
pathologies include provenance-free declaration, the Doll
Problem (declining exercise), and wireheading (declining broad
engagement); these share the structural-avoidance shape but
differ in formal detail.

The intended resolution is a \emph{proof-burden taxonomy} with
claim-and-challenge back-pressure.  Three burden types:
proof-of-non-mergeability (stuffing-specific),
proof-of-provenance (blocking pure-label declarations), and
proof-of-use (blocking persistently-dormant declarations
outside the residual class).  Proof-burdens operate as
canonicalization preprocessing on $P$ before M5 is applied;
M5's formula is preserved, but the preprocessing introduces
substantive new axiomatic content.  Claims are temporal and
contestable: a prior proof-of-non-mergeability can be voided by
subsequent discovery of a merger mechanism, with the discoverer
receiving structural-discovery credit
per~\citep[Proposition~2]{lasser2026horizon} and the avoider
losing the M5 credit the voided proof supported.  This makes
structural avoidance economically risky rather than
economically stable.

The sequence provides substrate but not implementation.
Paper~5 provides commitment infrastructure; the companion
paper~\citep{lasser2026paper8a} provides partial use-evidence
via the sacrifice aggregate (incomplete, because the
persistent-dormancy attack and private time-sacrifice
invisibility remain); Paper~6 provides a controlled-relaxation
pattern for adjudication.  A dedicated follow-up paper must
deliver: the formal canonicalization operator, the three
proof-burden types with admissible evidence classes, the
claim-and-challenge protocol with challenge-market robustness
(admission filtering against spam, identity weighting against
Sybil flooding, independence requirements against cartel
suppression, bounds on governance capture of the adjudication
layer), the discovery-reward calibration integrating with
Paper~3's structural-discovery result, and the
structural-avoidance dominance theorem under explicit
quantitative conditions (minimum challenger arrival rate,
minimum reward-to-cost ratio, maximum adjudication lag,
challenger independence).  Honest residuals include
information-asymmetry persistence, private time-sacrifice
invisibility inherited from the companion paper's P5 scope,
collusive-subsidy and external-funding edge cases,
canonicalization order/non-uniqueness, and the
deceptive-alignment residual that proof-burdens cannot
currently address without interpretability.  This is
substantial invention, grounded in existing sequence substrate
but not reducible to it; the natural venue is a dedicated
follow-up paper on the structural-avoidance architecture.

\item \label{oq:gap_decomp_closure}
\textbf{Automated verification of poset-boundary closure.}
The default priority-extension strategy
(Remark~\ref{rem:cell_cooperative_closure}(a)) produces
cooperative- and subsumption-closed cells by construction, but
verifying closure on adversarially-structured posets may be
non-trivial.  A structural-verification protocol
(analogous to controlled-relaxation testing against poset
claims) is a natural follow-up.

\item \label{oq:oi_bridge}
\textbf{OI--exercise bridge conditions.}
Proposition~\ref{prop:oi_floor} transfers Paper~3's
observational-individuation floor to $\volRexact^{[W]}(O)$
conditionally on the OI--exercise bridge
(Definition~\ref{def:oi_exercise_bridge}).  Where the bridge
fails (dormant presence), the structural $\volL$ floor still
holds but $\volRexact$ receives no contribution.  A cleaner
characterization of the bridge's domain of validity, and
conditions under which Paper~3's floor transfers directly
without the bridge, is open.

\end{enumerate}
